\documentclass[12pt,a4paper]{article}
\usepackage[utf8]{inputenc}
\usepackage{amsmath, amssymb, amsfonts}
\usepackage{geometry}
\geometry{margin=2.5cm}
\usepackage{hyperref}
\usepackage{pgfplots}
\usepackage{fancyhdr}
\usepackage{listing}
\usepackage{listings}
\usepackage{titling}
\usepackage{float}
\usepackage{titlesec}

\titleformat{\section}
  {\normalfont\Large\bfseries\centering}  % Format: font size, bold, centered
  {\thesection}{1em}{}
\pagestyle{fancy}
\pgfplotsset{compat=1.18}
\fancyhead{}
\fancyfoot{}
\title{Parlare in pubblico}
\pretitle{\begin{center}\LARGE}
\posttitle{\par\large Career Labs\end{center}}
\author{}
\date{}
\fancyhead[L]{Gabriele Perna}          % Left side
\fancyhead[C]{Parlare in pubblico}     % Center
\fancyhead[R]{2026/2027}             % Right side
\fancyfoot[C]{\thepage}           % Page number in the center of footer
\let\oldsection\section
\renewcommand{\section}{\clearpage\oldsection}

\begin{document}

\maketitle
\tableofcontents

\section{Informazioni importanti}
Ecco i contatti della relatrice:\\
\textbf{Email:} antonella.magliocchi@unipi.\\
\textbf{Linkedin:} Antonella Magliocchi

\section{Come mai si parla?}
Si parla per vari motivi:
\begin{itemize}
    \item Farsi capire
    \item Essere ricordati
    \item Generare un'azione
\end{itemize}

\section{Le funzioni della comunicazione (modello di Jakobson 1956)}
Il modello specifica queste entità:
\begin{itemize}
    \item \textbf{Emittente:} chi manda il messaggio
    \item \textbf{Messaggio:} contenuto della comunicazione
    \item \textbf{Destinatario:} chi riceve il messaggio
    \item \textbf{Contesto:} la situazione in cui avviene la situazione, bisogna conoscerlo per essere efficaci nel comunicare dei concetti
    \item \textbf{Referente:} il tema di cui si parla
    \item \textbf{Canale:} il mezzo attraverso cui passa la comunicazione (aula, meeting online, palco...), influisce su come comunichiamo nei gesti e le movenze
    \item \textbf{Codice:} il linguaggio utilizzato, varia in base alle altre entità
\end{itemize}

\section{La struttura che funziona sempre}
Un intervento efficace ha tre elementi fondamentali:
\begin{itemize}
    \item \textbf{Apertura:} un "hook" è importante, la più difficile, deve interessare i destinatari, altrimenti li perdiamo subito. Da un motivo al pubblico per ascoltarti.
    \item \textbf{Idee chiave:} il cervello tende a ricordare pochi concetti, per ogni punto: 1 idea e 1 esempio.
    \item \textbf{Chiusura:} Cosa voglio che resti? Sintetizza il messaggio chiave e invita all'azione.
\end{itemize}
Un'errore comune è cercare di dire tutto, in quanto il tempo è limitato e troppo contenuto può avere meno intatto.\\
Un discorso deve essere semplice e chiaro, gli esperti comunicatori sono in grado di esprimere concetti difficili in modo semplice e accessibile a tutti.\\
Inoltre, è necessario assecondare il \textbf{tempo} che abbiamo a disposizione per esporre il nostro discorso, non dobbiamo preparare più materiale / concetti di quelli che possiamo esporre in uno specifico lasso di tempo.

\section{Cosa fare prima di parlare}
Se ci invitano a parlare, dobbiamo prepararci bene:
\begin{itemize}
    \item Arrivare in tempo
    \item Provare le apparecchiature in anticipo
    \item Mangiare in anticipo
    \item Parlare con qualcuno del pubblico per ridurre le distanze (anche qualcuno che conosciamo)
    \item Vestirsi in modo da sentirsi a proprio agio
\end{itemize}
\section{L'importanza dell'interpretazione}
Non importa solo \textbf{cosa} dici, ma anche \textbf{come} lo dici.\\
Vediamo un video importante che ci aiuterà a capirlo: \href{https://www.youtube.com/watch?v=69EMd4csZRY}{https://www.youtube.com/watch?v=69EMd4csZRY}\\
Possiamo notare che Barack Obama fa queste cose durante il suo discorso:
\begin{itemize}
    \item Non guarda in un punto fisso, avendo una platea enorme, cerca di abbracciare con lo sguardo tutti.
    \item Fa molte pause, fondamentali per respirare e codificare i messaggi in modo giusto.
    \item Fa molti gesti, da colore ai messaggi
    \item Utilizza l'ironia per connettersi meglio con il pubblico, in particolare come hook.
\end{itemize}
Non solo questo aiuta un discorso, ma addirittura il contrario lo distrugge:
\begin{itemize}
    \item Una voce piatta spegne l'atmosfera
    \item L'assenza di pause ci fa sbagliare
    \item Gesticolare poco e muoversi poco ci fa sembrare spenti
\end{itemize}

\section{Il microfono}
Se utilizziamo un microfono, bisogna tenere di conto queste cose:
\begin{itemize}
    \item Non va tenuto troppo vicino: rimbomba
    \item Non va tenuto troppo lontano: non si seguente
    \item Bisogna stare lontano dalle casse, si crea un ritorno
\end{itemize}

\section{Cos'altro conta?}
Ecco altri elementi a cui bisogna portare attenzione quando parliamo in pubblico:
\begin{itemize}
    \item Dress-code appropriato al contesto
    \item Contatto visivo nelle conversazioni 1-1
    \item Evitare di essere troppo eccessivi nei gesti (a meno che non ci siano condizioni mediche o cose simili).
    \item Tono di voce adatto al contesto
\end{itemize}

\section{Gestire l'ansia}
L'ansia è un elemento negativo che può da sola distruggere un'intera performance, per combatterla ci sono delle varie tecniche:
\begin{itemize}
    \item Memorizzare la struttura
    \item Respirare lentamente prima di iniziare
    \item Avere una postura aperta e rilassata
    \item Accettare di poter sbagliare, capita a tutti
\end{itemize}


\end{document}